\section{Introduction}
\label{ch:introduction}

The dominant paradigm behind modern language models is \textbf{autoregressive generation}: a model factorizes the probability of a sequence into a product of next-token conditionals and produces text one token at a time, from left to right \citep{radford2019language, brown2020language}. This factorization is what makes these models trainable at scale and fast at inference, and it is the reason they work as well as they do. It is also, as this thesis argues, the reason one popular approach to controlling them does not.

That same factorization is the source of a practical limitation. Because generation proceeds strictly left to right and each token is fixed the moment it is emitted, the model has no mechanism for revising an earlier decision in light of a later one: if the second half of a sentence reveals that the first half has drifted in logic or tone, the model cannot go back, only continue or discard and restart. This is no obstacle when the goal is a fluent continuation, which is what these models are trained to produce, but it becomes one the moment we want to \emph{control} what they produce.

The reason control is hard has a precise structure. Many of the properties we care about most, the overall sentiment of a passage, its topical focus, its adherence to a factual constraint, the absence of toxic content, are \textbf{global properties} of a completed sequence rather than local properties of any single token. A greedy, left-to-right decision about the next token cannot see such a property, because the object that would be evaluated against it, the finished text, does not yet exist when the decision is made. Enforcing a global constraint autoregressively therefore requires generating many complete sequences and keeping those that happen to satisfy it, which is wasteful when the constraint is easy to satisfy and effectively impossible when it is not. Steering a language model toward such global objectives while keeping its output fluent is the problem of \textbf{controllable text generation}, and it is the subject of this thesis.

\subsection{Established Routes to Control Generation and Their Costs}
\label{sec:intro-routes}

There are, broadly, two established ways of imposing control on a language model, and each carries a cost that motivates the search for a third.

The first is to bake the desired behaviour into the parameters, through supervised fine-tuning on curated data or, more commonly today, reinforcement learning from human feedback \citep{ouyang2022training}. This is effective and now standard, and the aligned assistants in wide use are its product, but the control it provides is \emph{static}: a model tuned to be helpful and non-toxic embodies those constraints and no others, and a genuinely new constraint requires new preference data and a new and expensive training run.

The second is to leave the model untouched and filter or re-rank its outputs, generating many independent candidates and discarding those that violate the constraint. This preserves the model but inherits the intractability described above, since for a constraint a fluent generation satisfies only rarely by chance, the number of candidates one must draw grows without practical bound.

What both approaches lack, and what would be genuinely valuable, is the ability to introduce a fresh constraint at the moment of generation, cheaply, without touching the underlying model and without generating a haystack to find a needle. This is the capability that the energy-based approach promises, and it is the promise that this thesis sets out to examine.

\subsection{Energy-Based Control and Its Central Assumption}
\label{sec:intro-ebm}

An elegant alternative reframes decoding itself. Rather than generating a sequence token by token, one can define a probability distribution over the space of all possible sequences and \emph{sample} from it directly. Following the standard formulation of \textbf{energy-based models} \citep{lecun2006tutorial, hinton2002training}, the distribution is written in the Boltzmann form
\begin{equation}
    p(\bm{x}) = \frac{1}{Z}\exp\bigl(-\energy(\bm{x})\bigr),
    \label{eq:intro-ebm}
\end{equation}
where $\energy(\bm{x})$ is the \textbf{energy} of a sequence $\bm{x}$, a scalar that is low for desirable sequences and high for undesirable ones, and $Z$ is a normalizing constant. The constant $Z$ is intractable to compute, because it sums over every possible sequence, but the great advantage of the sampling framework is that it need never be computed: the methods used to draw samples from \eqref{eq:intro-ebm} depend only on \emph{ratios} of probabilities or on gradients of the log-probability, in both of which the constant cancels.

The attraction of this view for controllable generation is that the energy can be assembled additively from independent pieces. One writes
\begin{equation}
    \energy(\bm{x}) = -\log p_{\text{LM}}(\bm{x}) + \lambda\, C(\bm{x}),
    \label{eq:intro-energy}
\end{equation}
where the first term is the negative log-likelihood of a frozen, pretrained language model and supplies \emph{fluency}, pulling samples toward text the model considers natural, and the second term $C(\bm{x})$ is an arbitrary differentiable \emph{constraint}, weighted by a coefficient $\lambda$ that trades off fluency against constraint satisfaction. Because the constraint is simply added on, it can be specified at the moment of generation and swapped freely. A practitioner who wants positive sentiment today writes $C$ as the negative log-probability of the positive class under a sentiment classifier; one who wants factual consistency tomorrow writes a different $C$; and in both cases the language model is untouched. This is the promise of plug-and-play, inference-time controllable generation, and it is an appealing one. A substantial and growing body of recent work pursues it, through gradient-guided decoders such as PPLM \citep{dathathri2020plug}, FUDGE \citep{yang2021fudge}, and GeDi \citep{krause2021gedi}, and, most directly related to this thesis, through sampling-based methods such as COLD \citep{qin2022cold} and MuCoLa \citep{kumar2022gradient}, which draw samples from the energy in \eqref{eq:intro-energy} using continuous relaxations and dynamics derived from Langevin sampling.

Every method in this family rests on one assumption, which is worth isolating. To sample from the energy in \eqref{eq:intro-energy} using a gradient-guided method, two things must be true. First, one must be able to compute a gradient of $-\log p_{\text{LM}}(\bm{x})$ that points, at least on average, in the direction of better sequences. Second, one must be able to \emph{follow} that gradient across the space of possible texts, moving from a worse sequence to a better one by steps the gradient recommends. Taken together, these amount to the presupposition that the frozen likelihood of an autoregressive language model, viewed as a function over whole sequences, defines a \emph{landscape that a sampler can actually navigate}: a surface with meaningful slopes, whose downhill direction can be computed locally and trusted to lead somewhere good. This presupposition is intuitive, and it is almost never tested directly in the literature that relies on it. Testing it, and explaining what is found, is the subject of this thesis. In the settings tested here that presupposition does not hold: following the input-embedding gradient proves no more effective than following a norm-matched random direction, the finding stated precisely in Section~\ref{sec:results-fallacy}. This outcome is stated here, before the research questions are posed, because the thesis is organized as an explanation of that result rather than as a reveal.

That explanation has a definite shape, and setting it out now is what makes the order of the experiments legible. A null of this kind admits three candidate causes, and they are not equivalent. Under \textbf{Hypothesis A} the target is at fault: the frozen likelihood, evaluated on the counterfactual sequences an in-place revision must consider, is not an object any local search could profit from, in which case the choice of proposal mechanism is beside the point. Under \textbf{Hypothesis B} the target is serviceable but the model does not contain what a local proposal needs, because teacher forcing shapes next-token conditionals on correct prefixes and never asks the joint likelihood to rank substitutions; on this reading the deficiency is one of training objective, and only a differently trained model could repair it. Under \textbf{Hypothesis C} the information is present in the frozen model and the parameterization discards it: differentiating with respect to the input embedding registers only part of what a substitution changes, while a proposal that reads the model's output side recovers the discarded part, and one that may additionally condition on the right context recovers more still. The results chapter is arranged as an elimination in that order. Gradient-free rescoring and Gibbs sampling on the identical energy dispose of A; re-analysing the same measurements under a derivative taken in the token-indicator coordinates, and substituting that surrogate into the same sampler on the same frozen weights, disposes of B; and a ladder of proposals run inside one exact-energy chain confirms and refines C.

\subsection{From Sampling to Amortization}
\label{sec:intro-amortize}

The natural response to a gradient that will not cooperate is to stop using the gradient. There is a version of the energy-based idea that does not differentiate the energy at all. If a per-sequence energy can still be \emph{evaluated}, that is, if one can take a finished sequence and compute its score, even when one cannot usefully \emph{differentiate} that score to decide how to improve a sequence, then one can treat the language model as a black box that ranks finished sequences and learn a separate policy that generates in proportion to that ranking. This is the idea of \textbf{amortized inference}: rather than paying the cost of a search procedure afresh for every sample, one pays a one-time training cost to learn a generator that produces good samples directly.

The most principled recent realization of this idea for sequence generation is the \textbf{Generative Flow Network}, or GFlowNet \citep{bengio2021flow, bengio2023gflownet}, which trains a generative policy so that the probability of producing a complete sequence is proportional to its reward. Applied to language models, GFlowNets have been proposed as a way to sample diverse, high-reward text without the mode-collapse that afflicts reward-maximizing reinforcement learning \citep{hu2024amortizing}. The second half of this thesis asks a specific question about this approach: does amortizing the search, and thereby never touching the gradient, escape the difficulties that defeat the gradient-guided samplers? The answer, established in Section~\ref{ch:results}, is that it does not, and it separates into a statement about the learned policy and a statement about the energy that policy leaves behind. Because the initial sampling experiments failed, the study developed into a diagnostic investigation of the underlying mechanisms.

\subsection{Research Questions and Contributions}
\label{sec:intro-rqs}

The work is organized around three research questions, the third of which separates into two that the earlier drafts of this study conflated. The results section answers them in turn and the discussion section revisits each explicitly.

\begin{description}
    \item[RQ1.] Does the input-embedding gradient of a frozen autoregressive sequence likelihood provide a useful proposal direction for discrete token revision, and if not, which alternative local quantities do?
    \item[RQ2.] What is the role of the Metropolis--Hastings correction in each of the discrete and continuous samplers, and does equipping them with the exact accept-reject correction for the proposal they implement also make them work?
    \item[RQ3a.] Does a GFlowNet trained on the frozen likelihood as its reward learn a policy that generates high-reward text?
    \item[RQ3b.] Does the energy induced by that tuning become more navigable by the local input-embedding surrogate than the energy of the untuned base?
\end{description}

RQ3 is split because a GFlowNet is an amortized \emph{policy} while the Langevin samplers navigate an \emph{energy}, and running the samplers on a tuned model tests the geometry of that tuned energy rather than whether the policy amortizes sampling successfully. The two are related but not identical, and reporting them as one answer would blur a distinction the evidence supports.

One further question belongs to the study but not to its central diagnosis, and it is posed here as an extension.

\begin{description}
    \item[E.] Can an additive constraint term steer generation toward a target attribute on this energy landscape, as the plug-and-play framework requires?
\end{description}

This is the capability the whole programme was meant to deliver, and the thesis reports what happened when it was attempted; but it is an application of the sampling machinery whose prerequisite RQ1 tests, so it is answered after the primary diagnosis rather than alongside it, and the contributions below do not rest on it.\footnote{These questions are the operationalized forms of those posed in the thesis proposal, refined once the study became diagnostic rather than constructive. The proposal's RQ1 and RQ1a, on whether faithful Langevin dynamics can sample the frozen likelihood and how the two samplers compare, are consolidated into RQ1; its RQ1b, on the Metropolis--Hastings correction, becomes RQ2; its RQ3a, on amortizing the energy with a GFlowNet, becomes RQ3a and RQ3b; its RQ2a, on additive constraint steering, becomes the extension question E; and its RQ2b, on whether energy-based recovery improves on ordinary autoregressive decoding, is answered by the gradient-free baselines of Section~\ref{sec:results-baselines}.}

In answering these questions the thesis makes five contributions.

\begin{enumerate}
    \item A controlled ablation separating the \emph{direction} of the input-embedding gradient from its \emph{magnitude}, run across five energy functions and extended by a sweep over step size and proposal temperature that spans proposal entropies from uniform to deterministic. The direction gave no advantage over a norm-matched random one anywhere in that range, an equivalence certified at a margin fixed in advance over a thousand paired sequences.
    \item An identification of what the input-embedding derivative discards, and a demonstration that it is recoverable. The first-order surrogate is uncorrelated with the true energy change over the entire admissible range of substitution distances, and the self-and-future decomposition shows why; re-analysing the identical measurements under a derivative taken in the token-indicator coordinates restores a strong rank correlation, and substituting that surrogate into the same sampler on the same frozen weights lifts exact recovery from at most two percent to forty.
    \item A decomposition of the Metropolis--Hastings acceptance ratio in both samplers, which attributes the continuous sampler's collapse to the proposal term rather than the target term and shows that in the discrete sampler the accept/reject filter, not the gradient, supplies the entire selection pressure. Alongside it, a measurement of the \emph{likelihood trap} on the study's own models, in which the lowest-energy text is the most degenerate.
    \item A documented taxonomy of three failure modes of GFlowNet fine-tuning on a frozen-likelihood reward, together with the separate energy-level experiment showing that Langevin sampling on the amortized energy is indistinguishable from sampling on the untuned base even though tuning moved that energy by tens of nats.
    \item A controlled proposal ladder inside one exact-energy chain, in which recovery rises with what the proposal may condition on rather than with how it was trained: from the input-embedding gradient and a uniform draw at zero, through the frozen model's own left-to-right conditional, to a score-trained diffusion model and a masked language model that was never score-trained and performs best of all.
\end{enumerate}

\subsection{Structure of the Thesis}
\label{sec:intro-structure}

Section~\ref{ch:background} develops the technical background: the energy-based view of generation, Langevin dynamics, the Metropolis--Hastings correction, the two samplers studied here, discrete diffusion, and GFlowNets. Section~\ref{ch:related} reviews the literature on controllable and energy-based text generation and states the gap this work addresses. Section~\ref{ch:methodology} describes the task, the five energy functions, the sampler configurations, the evaluation metrics, and the diagnostic experiments. Section~\ref{ch:results} follows the elimination just described: it calibrates the samplers and establishes the null, disposes of Hypothesis A with gradient-free baselines on the identical energy, disposes of Hypothesis B with the token-indicator re-analysis and the sampler built on it, and closes with the proposal ladder that confirms Hypothesis C, alongside the GFlowNet experiments and the constrained-generation extension. Section~\ref{ch:discussion} answers each research question, sets out the shared problem and the distinct mechanisms beneath it, and states the limitations and implications. Section~\ref{ch:conclusion} concludes.
